\begin{lemma}[Well-founded prefix tree]
For every front $F$ on $X$, let
$T(F)=\{s\mid s\in\mathsf{PrefixTree}(F)\}$, regarded as a subtype of
finite subsets of $\mathbb N$.  The relation on $T(F)$ given by
\[
  r(u,v)\quad\Longleftrightarrow\quad
  \mathsf{ProperInitialSegment}(v,u)
\]
is well-founded.  Thus a recursive call from a node $v$ can proceed to a
strict extension $u$ of $v$; equivalently, $T(F)$ has no infinite branch.
\end{lemma}
\begin{proof}
Suppose an infinite branch existed.  Write its successive finite nodes as
$s_0\subsetneq s_1\subsetneq\cdots$, and let $X'$ be the infinite set formed
by their successive new elements.  Each $s_i$ is an initial segment of $X'$.
By the definition of \noderef{PrefixTree}, every such node is an initial
segment of a member of $F$, so $X'\subseteq\mathsf{FrontBase}(F)$.

If $F$ is trivial, its prefix tree contains only $\emptyset$, contradicting
the branch at length one.  Otherwise the base clause of \noderef{Front} gives
$X'\subseteq X$.  Density in \noderef{Front} supplies $s\in F$ with
 $s\segs X'$.  Choose a branch node $u$ with more elements than $s$.
 Since both are initial segments of $X'$, $s$ is a proper initial segment of
 $u$.  Since $u\in T(F)$,
\noderef{PrefixTree} supplies $t\in F$ with $u$ an initial segment of $t$.
Thus $s$ is an initial segment of $t$.  Prefix-freeness in
\noderef{Front}, together with $s,t\in F$, gives $s=t$, while
$s\segs u\segm t$ gives $s\ne t$, a contradiction.  Hence there is no
infinite branch, and the proper-extension relation is well-founded.
\end{proof}
